%% ============================================================
%% CV - Kossi Samuel Gabiam
%% Candidature - Programme de Recrutement des Jeunes Professionnels BIDC 2026
%% Compile: pdflatex cv_gabiam_fr.tex
%% ============================================================

\documentclass[10pt,a4paper]{article}

% ============================================================
% PACKAGES
% ============================================================

\usepackage[T1]{fontenc}
\usepackage[utf8]{inputenc}
\usepackage[french]{babel}
\usepackage{mathptmx}

\usepackage[
    margin=1.0cm,
    top=0.9cm,
    bottom=0.9cm
]{geometry}

\usepackage{hyperref}
\usepackage{enumitem}
\usepackage{titlesec}
\usepackage{tabularx}
\usepackage{array}
\usepackage{xcolor}
\usepackage{parskip}
\usepackage{fancyhdr}
\usepackage{needspace}

% ============================================================
% COULEURS
% ============================================================

\definecolor{sectioncolor}{HTML}{1E40AF}
\definecolor{secondary}{HTML}{475569}
\definecolor{linkcolor}{HTML}{1E40AF}

% ============================================================
% HYPERLIENS
% ============================================================

\hypersetup{
    colorlinks=true,
    urlcolor=linkcolor,
    linkcolor=linkcolor,
    pdfborder={0 0 0}
}

% ============================================================
% PAGES
% ============================================================

\pagestyle{fancy}
\fancyhf{}
\renewcommand{\headrulewidth}{0pt}
\fancyfoot[C]{\small\color{secondary}\thepage}

% ============================================================
% SECTIONS
% ============================================================

\titleformat{\section}
    {\normalsize\bfseries\color{sectioncolor}\scshape}
    {}{0em}
    {}
    [\vspace{-4pt}\color{sectioncolor}\rule{\linewidth}{1pt}]

\titlespacing{\section}
    {0pt}
    {8pt}
    {4pt}

% ============================================================
% LISTES
% ============================================================

\setlist[itemize]{
    leftmargin=1.25em,
    itemsep=1.5pt,
    topsep=0.5pt,
    parsep=0pt,
    partopsep=0pt,
    label=\small\textbullet
}

% ============================================================
% MACRO ENTRÉE CV
% ============================================================

% #1 = intitulé
% #2 = période
% #3 = organisme
% #4 = localisation / type

\newcommand{\cventry}[4]{%
    \needspace{4\baselineskip}%
    \noindent
    \begin{tabularx}{\linewidth}{
        @{}
        X
        >{\raggedleft\arraybackslash}p{3.15cm}
        @{}
    }
        \textbf{\large #1}
        &
        {\small\color{secondary}\textit{#2}}
    \end{tabularx}
    \par
    \ifx&#3&
    \else
        \noindent
        \textbf{\textit{#3}}
        \par
    \fi
    \ifx&#4&
    \else
        \noindent
        {\small #4}
        \par
    \fi
    \vspace{2.5pt}%
}

% ============================================================
% DOCUMENT
% ============================================================

\begin{document}

\color{black}

% ============================================================
% EN-TÊTE
% ============================================================

\begin{center}

    {\LARGE\bfseries\scshape GABIAM Kossi Samuel}

    \par
    \vspace{4pt}

    {\normalsize\bfseries
    Systèmes, Réseaux \& Cloud
    \textbullet{}
    Développeur Full-Stack}

    \par
    \vspace{6pt}

    {\small
    \href{mailto:gabiam.k.samuel@gmail.com}
    {gabiam.k.samuel@gmail.com}
    \quad\textbullet\quad
    +228 90 84 77 68
    \quad\textbullet\quad
    Lomé, Togo
    }

    \par
    \vspace{2pt}

    {\small
    \href{https://www.linkedin.com/in/samuel-g-107663165/}
    {linkedin.com/in/samuel-g-107663165}
    \quad\textbullet\quad
    \href{https://github.com/Shieldish}
    {github.com/Shieldish}
    }

    \par
    \vspace{2pt}

    {\small
    \href{https://gitlab.com/Shieldish}
    {gitlab.com/Shieldish}
    \quad\textbullet\quad
    \href{https://portfolio.gabiam-k-samuel.workers.dev/}
    {portfolio.gabiam-k-samuel.workers.dev}
    }

\end{center}

\vspace{-2pt}

\noindent
\rule{\linewidth}{1pt}

\vspace{3pt}

% ============================================================
% PROFIL PROFESSIONNEL
% ============================================================

\section{Profil Professionnel}

Titulaire d'un
\textbf{Master Professionnel en Systèmes, Réseaux \& Cloud Computing}
et d'une
\textbf{Licence en Génie Logiciel},
je suis développeur Full-Stack spécialisé dans la conception
et l'intégration de systèmes informatiques, d'applications web
et mobiles et d'infrastructures distribuées.

Mon expérience couvre notamment les
\textbf{architectures distribuées, les API, les bases de données,
la sécurité applicative, l'authentification, le Cloud,
la conteneurisation et les pipelines DevOps}.
Je possède une approche orientée résolution de problèmes,
fiabilité des systèmes et amélioration des performances,
ainsi qu'une expérience dans le développement de solutions
utilisées en production au Togo.

% ============================================================
% EXPÉRIENCE PROFESSIONNELLE
% ============================================================

\section{Expérience Professionnelle}

\cventry
{Développeur Full-Stack}
{Juin 2025 -- Présent}
{SuiSco Sarl}
{Lomé, Togo}

\begin{itemize}

    \item Conception et développement de solutions web et de services
    numériques destinés à différents usages et canaux, avec une attention
    particulière portée à la \textbf{fiabilité, la sécurité et la performance}.

    \item \textbf{EduBoost :} architecture et développement d'une plateforme
    nationale multi-canaux intégrant Web, SMS, USSD et Mobile Money,
    actuellement en production au Togo.

    \item Conception d'une \textbf{architecture distribuée} composée de
    15 nœuds Django d'écriture et d'un nœud central d'agrégation,
    avec mécanisme de génération de tickets anti-collision basé sur PostgreSQL.

    \item Mise en œuvre de traitements \textbf{asynchrones et distribués}
    avec Celery, Redis et PostgreSQL, incluant synchronisation,
    reprise après incident et traitement des demandes Web.

    \item Intégration de mécanismes de \textbf{sécurité applicative} :
    OTP, HMAC-SHA256 pour la vérification des webhooks,
    rate-limiting Redis et authentification centralisée.

    \item Développement et maintenance de plusieurs solutions complémentaires :
    \textbf{DefiQuiz}, SuiSco My Account et intégration du
    \textbf{Single Sign-On Keycloak} basé sur OAuth2/OIDC.

\end{itemize}

% ============================================================
% PROJETS MAJEURS
% ============================================================

\section{Projets Majeurs}

\cventry
{DefiQuiz -- Plateforme de Quiz Éducatif}
{2025 -- Présent}
{SuiSco Sarl / YAS TOGO}
{Plateforme Web}

\begin{itemize}

    \item Développement d'une plateforme de quiz multi-clients
    proposant plus de \textbf{5\,000 questions}.

    \item Conception du frontend et du backend avec
    Nuxt.js / Next.js, Hono.js et Deno.js.

\end{itemize}

\cventry
{SuiSco My Account \& SSO Centralisé}
{2025}
{SuiSco Sarl}
{Applications internes}

\begin{itemize}

    \item Développement d'une application web multi-clients de gestion
    de comptes en PHP Laravel, API REST et MySQL.

    \item Mise en place d'un système de \textbf{Single Sign-On}
    avec Keycloak et OAuth2/OIDC pour centraliser et sécuriser
    l'authentification.

\end{itemize}

\cventry
{Plateforme de Gestion de Stages}
{2024}
{Djagora}
{Projet de Fin d'Études}

\begin{itemize}

    \item Développement d'une application Web et Mobile pour la gestion
    des offres de stage, des candidatures et du suivi des étudiants.

    \item \textbf{Technologies :} React, Node.js, MySQL et React Native.

\end{itemize}

\cventry
{Extraction \& Visualisation de Données Sociales}
{2021}
{DK-Soft}
{Projet de Stage}

\begin{itemize}

    \item Développement d'une application d'extraction et de visualisation
    analytique en temps réel de données issues des réseaux sociaux.

    \item Mise en œuvre de la \textbf{Stack ELK} :
    Elasticsearch, Logstash et Kibana, avec Angular.

\end{itemize}

% ============================================================
% FORMATION ACADÉMIQUE
% ============================================================

\section{Formation Académique}

\cventry
{Master Professionnel en Systèmes, Réseaux \& Cloud Computing}
{2022 -- 2024}
{Faculté des Sciences de Sfax}
{Sfax, Tunisie}

\begin{itemize}

    \item Ingénierie des systèmes informatiques, architectures distribuées,
    virtualisation, Cloud Computing, administration Linux et sécurité réseau.

\end{itemize}

\cventry
{Licence Nationale en Science de l'Informatique}
{2019 -- 2022}
{Faculté des Sciences de Sfax}
{Sfax, Tunisie}

\begin{itemize}

    \item Spécialité \textbf{Génie Logiciel \& Systèmes d'Information}.

    \item Algorithmique, bases de données, UML, génie logiciel
    et réseaux informatiques.

\end{itemize}

\cventry
{Baccalauréat Scientifique -- Série D}
{2014 -- 2017}
{Lycée La Grâce}
{Sokodé, Togo}

% ============================================================
% COMPÉTENCES TECHNIQUES
% ============================================================

\section{Compétences Techniques}

\begin{tabularx}{\linewidth}{
    @{}
    >{\bfseries\raggedright\arraybackslash}p{3.0cm}
    @{\hspace{8pt}}
    >{\raggedright\arraybackslash}X
    @{}
}

Développement &
React, Next.js, Vue.js, Nuxt 3, Angular,
TypeScript, Java, Go, Django, DRF, Node.js,
PHP Laravel, Hono.js, Deno.js, Python,
API REST, API GraphQL
\\[3pt]

Bases de données &
PostgreSQL, MySQL, Redis, MongoDB,
Elasticsearch
\\[3pt]

Systèmes \& Cloud &
Linux/Unix, Docker, Docker Compose,
architectures distribuées, virtualisation,
Keycloak, SSO
\\[3pt]

DevOps \& Méthodes &
Git, GitHub, GitLab, CI/CD, Celery,
Gunicorn, Agile / Scrum
\\[3pt]

Sécurité &
OAuth2/OIDC, HMAC-SHA256, OTP,
rate-limiting, authentification et sécurisation
des API
\\[3pt]

Intelligence artificielle &
Outils d'IA générative pour le développement,
la recherche, l'analyse, la documentation et
la résolution de problèmes.
Capacité à concevoir, entraîner et
expérimenter des modèles d'IA.
\\[3pt]

Mobile \& Data &
React Native, Flutter, ELK, Elasticsearch,
Kibana, Loki, Grafana, Prometheus

\end{tabularx}

% ============================================================
% LANGUES
% ============================================================

\vspace{3pt}

\noindent
\textbf{Langues :}
Français (langue maternelle)
\quad\textbullet\quad
Anglais (compétence professionnelle)

% ============================================================
% FIN
% ============================================================

\end{document}